% ============================================================
% Results chapter — Contract-Year Effect in Professional Football
% Drop-in for Overleaf. Requires only: \usepackage{booktabs}.
% Self-contained; no siunitx needed.
% ============================================================

\section{Results}
\label{sec:results}

This chapter reports the empirical findings. It proceeds from the raw
descriptive contrast between players inside and outside the Bosman window,
through the intensive-margin analysis of match performance ratings, to the
extensive-margin analysis of playing time, and finally to the contract-cycle
event study that serves as the study's main exhibit. The chapter also documents
the designs that were tested and rejected---the regression discontinuity around
the six-month threshold and the naive contract-renewal specification---because
their failure is itself informative about the identification problem. The
central result can be stated at the outset: the data reveal a large, robust
decline in the \emph{quantity} of football a player is given to play as his
contract runs down, but no detectable change in per-minute \emph{performance}
once selection into playing is taken seriously. The incentive channel emphasised
in the problem statement, in other words, operates on playing time rather than
on measured effort.

\subsection{Descriptive patterns}
\label{sec:results-descriptive}

The estimation sample is a player-month panel of 152{,}817 observations
covering 5{,}695 players and 1{,}448 clubs across major European leagues,
spanning June 2022 to March 2026; 14{,}298 of these player-months fall inside
the six-month Bosman window. A first, unconditional comparison of means
(Table~\ref{tab:bosman-means}) already sketches the pattern that the remainder
of the chapter refines. Players observed inside the Bosman window are, on
average, 2.7 years older than those outside it, consistent with contract expiry
being reached later in a career. They play meaningfully less---roughly nineteen
fewer minutes per month---and register slightly fewer goals and assists.

The performance-rating contrast is where interpretation must already become
careful. The unweighted mean monthly rating is 0.14 points lower inside the
window and the difference is statistically significant, but the minutes-weighted
rating, which down-weights the short substitute appearances that dominate
low-rating observations, is essentially identical across the two groups
($+0.02$, $p = 0.124$). The raw rating gap is therefore not obviously a
performance gap: it is confounded with playing time, with age, and with the
composition of who is observed playing at all. Disentangling these is the task
of the regression analysis.

\begin{table}[htbp]
  \centering
  \caption{Player-months inside vs.\ outside the Bosman window (unconditional means).}
  \label{tab:bosman-means}
  \begin{tabular}{lrrrr}
    \toprule
    Variable & Bosman & Non-Bosman & Diff. & $p$ \\
    \midrule
    Days to expiry        & 116.7 & 852.8 & $-736.1$ & ${<}0.001$ \\
    Age (years)           & 27.63 & 24.92 & $+2.71$  & ${<}0.001$ \\
    Minutes (month)       & 146.0 & 164.7 & $-18.7$  & ${<}0.001$ \\
    Goals (month)         & 0.21  & 0.26  & $-0.06$  & ${<}0.001$ \\
    Assists (month)       & 0.15  & 0.19  & $-0.04$  & ${<}0.001$ \\
    Mean rating           & 5.79  & 5.93  & $-0.14$  & ${<}0.001$ \\
    Minutes-wtd.\ rating  & 6.68  & 6.66  & $+0.02$  & $0.124$    \\
    \bottomrule
  \end{tabular}

  \vspace{4pt}
  {\footnotesize\itshape Note: strict all-competitions panel; two-sample
  $t$-tests. Ratings use positive observations only.\par}
\end{table}

\subsection{The intensive margin: match performance}
\label{sec:results-intensive}

The intensive-margin question is whether a player, when he does take the field,
performs better as contract expiry and the associated free-agency option
approach. The preferred outcome is the monthly match rating standardised within
league-season cells, which removes the cross-league scale and composition
differences that otherwise contaminate a pooled estimate; all models are
estimated with player and time fixed effects and standard errors clustered by
player.

Pooled across the sample, there is no Bosman performance effect. The point
estimate on the Bosman indicator is $+0.041$ standard deviations ($p = 0.06$) in
the all-competitions sample and a precisely estimated $+0.004$ ($p = 0.87$) when
attention is restricted to players' own domestic-league matches. Neither is
distinguishable from zero at conventional levels, and the sign of the marginal
all-competitions estimate is the opposite of what a within-player decline would
require.

A career-concerns reading would predict that any effort response is concentrated
among younger players, for whom the labour market's evaluation carries the
highest option value. The age-interacted estimates are, at first glance,
encouraging: the Bosman coefficient is $+0.198$ ($p = 0.011$) for players under
23 and $+0.128$ ($p < 0.001$) for those aged 23--28, while the over-29 group
shows nothing. These pooled positives do not survive the study's robustness
protocol, however. Re-estimating the under-23 interaction season by season, no
individual season is significant, and the estimate declines monotonically as
measurement coverage improves---from $+0.26$ in 2023/24, to $+0.18$ in 2024/25,
to $+0.06$ ($p = 0.61$) in 2025/26, the best-measured season. The pooled result
rests on only 533 under-23 Bosman player-months from 219 players and is best
read as a composition artefact of the panel's changing league mix rather than as
evidence of a behavioural response.

The specification grid reinforces this conclusion. When identification is
tightened from player fixed effects to player-by-contract-spell fixed
effects---so that the comparison is made within a single contract, along which
time-to-expiry advances mechanically and cannot be a chosen state---the Bosman
coefficient moves from $+0.041$ (player $+$ month) and $+0.047$ (player $+$
league-month) to $-0.020$ (spell $+$ month) and $-0.011$ (spell $+$
league-month), all statistically insignificant. The sign is unstable across
fixed-effects structures and every spell-based estimate is null. This is not an
underpowered null: the minimum detectable effect at 80\% power is approximately
0.06--0.07 standard deviations, so effects larger than that are ruled out with
good power. Two-way clustering by player and month, and standardisation within
league-season-position cells, leave the picture unchanged.

\subsection{What the rating measures, and why the null is credible}
\label{sec:results-anatomy}

Interpreting an absence of effect requires knowing what the outcome variable
actually captures. A decomposition of roughly 337{,}000 rated player-matches
is instructive. Regressing the match rating on its observable ingredients,
minutes played alone explain 15\% of the variance, goals and assists raise this
to 43\%, and adding the average rating of a player's teammates that day reaches
55\%; the implied prices are about $+0.93$ rating points for a goal, $+0.77$ for
an assist, $-1.03$ for a red card, and $+0.90$ for a full ninety minutes. In
round terms the rating is roughly 45\% scoreboard events and minutes, 10\% the
team's collective result, and 45\% residual. Three features of this structure
bear directly on the analysis.

First, a single match rating is mostly noise about the individual: the
within-player match-to-match standard deviation (0.74) is more than twice the
between-player standard deviation of career averages (0.31), so a monthly
average of about four matches is only around 40\% signal. Second, roughly a
third of match-level rating variance sits at the team-match level, because
results and match dominance enter every teammate's rating simultaneously; the
rating therefore partly measures team quality, which is why the strictest
specifications difference out team-time variation. Third, the rating is
effectively a different instrument in each position---explained variance ranges
from 0.75 for attackers, whose ratings are largely restated goal involvement,
down to 0.23 for goalkeepers, whose ratings are dominated by team
outcomes---and minutes buy rating points mechanically even in event-free
matches, with substitute appearances shrunk toward a common anchor near 6.2.

These properties mean the intensive-margin null is a statement about a
deliberately noisy, minutes-contaminated, team-laden measure. Correcting for the
contamination does not overturn it: outcomes built from sixty-minute-plus
appearances only, minutes-bin-rescaled ratings, and league-month percentile
ranks that are immune to skew and shrinkage all reproduce the same
signature---small positives under naive fixed effects that vanish under spell
fixed effects and are near-exact zeros within club-month. Correcting the
measurement sharpens the null rather than revealing a hidden effect.

\subsection{The extensive margin: playing time}
\label{sec:results-extensive}

Because a rating exists only when a player plays, and playing itself responds to
contract status, the intensive-margin regressions condition on an outcome of the
treatment. The natural response is to analyse the fully observed extensive
margin---whether a player plays at all, and how many minutes and matches he
accumulates---directly on the full panel, using the always-observed
Transfermarkt playing records. Here a single result survives every
fixed-effects structure. Players inside the Bosman window are 4.4 to 6.6
percentage points less likely to play in a given month, a reduction that remains
significant under the strictest specification with contract-spell and
league-month fixed effects ($-4.6$~pp, $p < 0.001$). The effect steepens sharply
in the final half-year: the zero-to-six-months-to-expiry bin carries an
18.8-percentage-point reduction in playing probability in the strictest
specification.

That the effect is a genuine within-team phenomenon, rather than a reflection of
which clubs happen to have expiring players, is confirmed by comparing a player
to his own teammates in the same month through club-by-month fixed effects. The
playing-probability penalty holds ($-2.4$~pp, $p < 0.001$) and minutes fall
($-5.3$, $p = 0.04$). The mean-rating outcome shows a marginal $-0.064$
($p = 0.04$) in the same specification, but this is the mechanical minutes
channel and not a performance effect: the minutes-weighted rating is null
($-0.017$, $p = 0.51$), and both are null among regulars playing 270 or more
minutes per month, where the minutes-rating gradient is flat.

The selection structure also disciplines the interpretation of the intensive
margin. Since expiring players are benched, the players who continue to appear
are positively selected, which biases any conditional rating comparison upward
relative to a true negative effect. A Lee-style bounding exercise makes the
limit explicit: trimming the over-observed control group by its differential
observation share within league-month cells produces bounds on the Bosman rating
gap of roughly $[-0.31,\, +0.19]$ standard deviations around an untrimmed
$-0.02$. The bounds comfortably span zero. Given the extensive-margin selection,
the conditional rating data cannot even sign a performance effect---which is
precisely why the extensive margin, and not the rating, is where the
identifiable action lies.

\subsection{The contract-cycle event study}
\label{sec:results-eventstudy}

The study's main exhibit estimates each outcome as a function of months
remaining to contract expiry, from twenty-four-plus down to zero with eighteen
months as the reference, within player-by-contract-spell and league-by-month
fixed effects. The playing-time profile is unambiguous. Playing probability and
minutes decline smoothly and monotonically beginning roughly twelve months
before expiry, reaching about $-25$ percentage points in playing probability and
$-70$ minutes per month in the expiry month itself. Crucially, there is no
discontinuous jump at the six-month Bosman threshold: the decline is a smooth
ramp through it. This is direct evidence for why a discontinuity design finds
nothing real at the legal cutoff---the theory-consistent pattern is anticipatory
and continuous, not discontinuous, because eligibility for free agency is
perfectly foreseen.

The standardised rating conditional on playing, by contrast, is flat across the
entire contract cycle, with confidence intervals straddling zero at every
horizon. The intensive-margin null is thus shown rather than merely asserted.
Two robustness checks guard against obvious confounds. The playing-time profile
is unchanged by quadratic age controls and is equally present for players who
were under 24 at the start of their contract spell---a group for whom advancing
age raises playing time---so the decline cannot be an aging artefact. And
although 88\% of expiry-month observations fall in June, restricting the sample
to active league-months leaves the profile essentially identical (month-zero
coefficient $-0.261$ versus $-0.255$), because the league-month fixed effects
already compare expiring players to same-league, same-month peers.

\subsection{Rejected designs}
\label{sec:results-rejected}

Two designs anticipated in the research plan were tested and are retained as
negative exhibits. The regression discontinuity at 180 days to expiry proves
unusable as causal evidence. Because the running variable is measured in monthly
steps, there is almost no support within a narrow bandwidth---120 observations
against 906 at a thirty-day bandwidth---and the headline local-linear estimate is
implausibly large (about $+19.3$ on a zero-to-ten scale, minutes-weighted), of
the opposite sign to the raw means gap it is meant to capture. Placebo cutoffs
confirm the artefact: a fake threshold at 150 days yields $+3.34$ ($p = 0.003$)
and one at 210 days $-2.09$ ($p = 0.002$). The estimates reflect local slope
extrapolation and composition around contract-time thresholds, not a
discontinuity at the legal Bosman cutoff, where by construction no discontinuous
treatment exists because eligibility is fully anticipated.

The naive contract-renewal specification fails a parallel audit. A pooled dummy
for the post-renewal period carries a significant $-0.097$ ($p < 0.001$), which
superficially suggests players relax after securing a new deal. Within every
individual season, however, the coefficient is insignificant with mixed signs,
and an event study around the first observed renewal shows within-league ratings
slightly \emph{rising} after signing (by about $+0.08$ standard deviations at
months four to six). The raw post-renewal average is likewise higher, not lower,
than the pre-renewal average (6.56 versus 6.17), reflecting selection---better
players get renewed---rather than a within-player decline. The pooled negative is
mean reversion combined with cross-season composition, and a staggered-adoption
estimator would be the appropriate formalisation; it is flagged as an extension
rather than presented as a finding.

\subsection{Financial heterogeneity}
\label{sec:results-financial}

A final set of tests asks whether the benching of expiring players varies with
club finances, using season-lagged wage and wage-to-revenue measures matched to
the panel. The benching pattern is confirmed within the financially documented
subsample---the reduction in playing probability is $-4.4$~pp ($p = 0.02$) in
high-wage clubs and $-4.8$~pp ($p = 0.02$) in low-financial-pressure
clubs---but no financial split is formally distinguishable, with joint
interaction tests at $p = 0.97$ for wage terciles and $p = 0.30$ for
wage-to-revenue terciles. A pre-declared continuous test offers a suggestive,
and only suggestive, refinement: the interaction of the Bosman indicator with a
club's wage-to-revenue percentile is $+0.070$ ($p = 0.051$), implying that the
benching is deepest (around $-7$~pp) at the most financially relaxed clubs and
shrinks as wage pressure rises. This is consistent with a squad-slack mechanism,
but at a Bonferroni-adjusted $p$ of roughly 0.10 it is reported as an avenue
rather than an established difference. An earlier significant wage-to-revenue
gradient obtained under improper look-ahead assignment did not survive correct
season-lagging and is retained only as a cautionary note on temporal assignment.

\subsection{Synthesis}
\label{sec:results-synthesis}

Taken together, the results describe an incentive that acts on opportunity
rather than on measured effort. As a contract runs down, a player is
progressively withdrawn from play---smoothly, from about a year out, and steeply
in the final months---and this extensive-margin decline is the one finding that
survives every specification, every fixed-effects structure, and comparison
against a player's own teammates. Per-minute performance, by contrast, shows no
detectable change; and given that the rating is a noisy, minutes-driven,
team-laden measure observed only for a positively selected set of players who
keep appearing, a performance effect could not have been credibly signed from
these data even had one been present. The tempting positive results---a
career-concerns gradient among the young, a post-renewal slackening, a
discontinuity at the legal threshold---each dissolve under season-stability,
specification, and placebo audits. The honest headline is therefore a
well-audited null on the effort margin paired with a robust effect on the
playing-time margin, with the caveat, applicable throughout, that the panel's
league composition shifts substantially across seasons and remains the single
largest threat to any pooled estimate.
